% !TeX program = pdflatex
\documentclass[11pt,a4paper]{article}
\usepackage[utf8]{inputenc}
\usepackage[T1]{fontenc}
\usepackage[english]{babel}
\usepackage{amsmath,amssymb,amsthm}
\usepackage[margin=2.4cm]{geometry}
\usepackage{booktabs}
\usepackage{array}
\usepackage{enumitem}
\usepackage{xcolor}
\usepackage[hidelinks,breaklinks]{hyperref}
\usepackage[expansion=false]{microtype}

\sloppy
\emergencystretch=2em

\newtheorem{definicao}{Definição}
\newtheorem{teorema}{Teorema}
\newtheorem{proposicao}{Proposição}
\newtheorem{lema}{Lema}
\newtheorem{observacao}{Observação}
\newtheorem{corolario}{Corolário}

\newcommand{\code}[1]{\texttt{\small #1}}
\newcommand{\abs}[1]{\lvert #1\rvert}

\title{Espaço de parâmetros algébricos\\
\large Família de álgebras quadráticas parametrizada por \(\mathbb{Q}\)}
\author{}
\date{}

\begin{document}
\maketitle

\section{Escopo e estatuto matemático}

Esta seção formaliza a família de parâmetros racionais associada aos polinômios
\[
P_q(\lambda)=\lambda^2-q\lambda-1,\qquad q\in\mathbb{Q},
\]
à sua realização matricial
\[
A_q=\begin{pmatrix}q&1\\1&0\end{pmatrix}
\]
e à álgebra gerada
\[
\mathcal{A}_q=\mathbb{Q}[A_q]\subset M_2(\mathbb{Q}).
\]

A postura metodológica é a seguinte: formalizar primeiro as estruturas que são
consequência direta das relações já estabelecidas; introduzir estruturas
adicionais somente quando forem necessárias e justificadas. Em particular,
não se definem operações artificiais entre matrizes de parâmetros distintos,
nem se afirma a existência de uma geometria canônica associada a cada \(q\).

A cadeia central é
\[
\boxed{
q\in\mathbb{Q}
\longmapsto
P_q
\longmapsto
A_q
\longmapsto
\mathcal{A}_q
}
\]
com as identificações
\[
\mathcal{A}_q
\cong
\mathbb{Q}[t]/(P_q)
\cong
\mathbb{Q}[w]/(w^2-\Delta_q),
\qquad
\Delta_q=q^2+4.
\]

A distinção fundamental é
\[
\boxed{
\text{dependência algébrica em }q
\;\neq\;
\text{dependência geométrica real em }q.
}
\]

\section{Família \(P_q\), \(q\in\mathbb{Q}\)}

\begin{definicao}
Para cada \(q\in\mathbb{Q}\) define-se o polinômio
\[
P_q(\lambda)=\lambda^2-q\lambda-1.
\]
O espaço de parâmetros é o conjunto
\[
\mathcal{P}_{\mathbb{Q}}
=
\bigl\{P_q:q\in\mathbb{Q}\bigr\}.
\]
\end{definicao}

A aplicação
\[
\Phi:\mathbb{Q}\longrightarrow\mathcal{P}_{\mathbb{Q}},
\qquad
\Phi(q)=P_q
\]
é bijetiva, pois o coeficiente de \(\lambda\) determina univocamente \(q\).
Assim,
\[
\mathcal{P}_{\mathbb{Q}}\cong\mathbb{Q}
\]
como conjuntos. A estrutura aqui é apenas a de parametrização; nenhuma
topologia, métrica ou estrutura diferencial é implícita.

No nível inteiro define-se analogamente
\[
\mathcal{P}_{\mathbb{Z}}=\{P_m:m\in\mathbb{Z}\}\subset\mathcal{P}_{\mathbb{Q}}.
\]

\section{Realização matricial \(A_q\)}

\begin{definicao}
Para cada \(q\in\mathbb{Q}\) define-se a matriz companheira
\[
A_q=\begin{pmatrix}q&1\\1&0\end{pmatrix}\in M_2(\mathbb{Q}).
\]
\end{definicao}

Tem-se imediatamente
\begin{align*}
\operatorname{tr}(A_q)&=q,\\
\det(A_q)&=-1,
\end{align*}
e o polinômio característico
\[
\chi_{A_q}(\lambda)=\det(\lambda I-A_q)=P_q(\lambda).
\]
Logo
\[
\boxed{\chi_{A_q}=P_q.}
\]

Pelo teorema de Cayley--Hamilton,
\[
A_q^2-qA_q-I=0,
\]
ou equivalentemente
\[
\boxed{A_q^2=qA_q+I.}
\]

A matriz
\[
T_q=\begin{pmatrix}0&1\\1&q\end{pmatrix}
\]
é semelhante a \(A_q\) via a permutação
\[
J=\begin{pmatrix}0&1\\1&0\end{pmatrix},
\qquad
T_q=JA_qJ.
\]
A escolha entre \(A_q\) e \(T_q\) é apenas de representação matricial.

\section{Álgebra \(\mathcal{A}_q\)}

\begin{definicao}
A álgebra gerada por \(A_q\) é o subespaço
\[
\mathcal{A}_q:=\operatorname{span}_{\mathbb{Q}}\{I,A_q\}\subset M_2(\mathbb{Q}).
\]
\end{definicao}

\begin{proposicao}
\(\mathcal{A}_q\) é uma \(\mathbb{Q}\)-álgebra comutativa de dimensão \(2\),
com elemento identidade \(I\). Explicitamente, todo elemento escreve-se de
modo único como \(\alpha I+\beta A_q\) com \(\alpha,\beta\in\mathbb{Q}\), e
a multiplicação é determinada pela relação \(A_q^2=qA_q+I\).
\end{proposicao}

\begin{proof}
A independência linear de \(I\) e \(A_q\) segue imediatamente da forma das
matrizes. O fechamento multiplicativo decorre da relação de Cayley--Hamilton.
A comutatividade é imediata.
\end{proof}

\begin{teorema}
O mapa
\[
\varphi:\mathbb{Q}[t]/(t^2-qt-1)\longrightarrow\mathcal{A}_q,
\qquad
\overline{t}\longmapsto A_q
\]
é um isomorfismo de \(\mathbb{Q}\)-álgebras.
\end{teorema}

\begin{proof}
O mapa é bem-definido porque \(A_q\) satisfaz \(P_q\). É sobrejetivo por
construção de \(\mathcal{A}_q\). É injetivo porque ambos os espaços têm
dimensão \(2\).
\end{proof}

Assim,
\[
\boxed{\mathcal{A}_q\cong\mathbb{Q}[t]/(P_q)}.
\]
Não se trata de uma álgebra nova: é a realização matricial canônica da
álgebra quociente associada ao polinômio \(P_q\).

As potências de \(A_q\) admitem a expressão
\[
A_q^n=a_n(q)A_q+b_n(q)I,
\]
onde as sequências \((a_n)\) e \((b_n)\) satisfazem a recorrência
\[
a_{n+2}=qa_{n+1}+a_n
\]
com condições iniciais adequadas.

\section{Apresentação equivalente via \(W_q\)}

\begin{definicao}
Define-se o operador
\[
W_q=2A_q-qI=\begin{pmatrix}q&2\\2&-q\end{pmatrix}.
\]
\end{definicao}

Tem-se a relação inversa
\[
A_q=\frac{W_q+qI}{2}.
\]

\begin{proposicao}
\[
W_q^2=(q^2+4)I.
\]
\end{proposicao}

\begin{proof}
Expandindo e usando \(A_q^2=qA_q+I\) obtém-se
\begin{align*}
W_q^2
&=(2A_q-qI)^2
=4A_q^2-4qA_q+q^2I
=4(qA_q+I)-4qA_q+q^2I
=(q^2+4)I.
\end{align*}
\end{proof}

Como \(\{I,W_q\}\) também gera \(\mathcal{A}_q\), tem-se
\[
\mathcal{A}_q=\operatorname{span}_{\mathbb{Q}}\{I,W_q\}.
\]

\begin{teorema}
O mapa
\[
\psi:\mathbb{Q}[w]/(w^2-(q^2+4))\longrightarrow\mathcal{A}_q,
\qquad
\overline{w}\longmapsto W_q
\]
é um isomorfismo de \(\mathbb{Q}\)-álgebras. Portanto
\[
\boxed{\mathcal{A}_q\cong\mathbb{Q}[w]/(w^2-\Delta_q)},
\qquad
\Delta_q=q^2+4.
\]
\end{teorema}

As duas apresentações são isomorfas via a mudança de variável \(w=2t-q\).

\section{Involução e paridade interna}

A relação \(A_q^2=qA_q+I\) determina uma involução natural na álgebra \(\mathcal{A}_q\).

\begin{definicao}
Define-se
\[
\overline{A_q}:=qI-A_q=\begin{pmatrix}0&-1\\-1&q\end{pmatrix}.
\]
\end{definicao}

\begin{proposicao}
Tem-se
\[
\overline{\overline{A_q}}=A_q
\qquad\text{e}\qquad
A_q\overline{A_q}=\overline{A_q}A_q=-I.
\]
Em particular \(\overline{A_q}=-A_q^{-1}\).
\end{proposicao}

\begin{proof}
A primeira identidade é imediata. A segunda segue de
\begin{align*}
A_q\overline{A_q}
&=A_q(qI-A_q)
=qA_q-A_q^2
=qA_q-(qA_q+I)
=-I.
\end{align*}
\end{proof}

Aplicando a mesma involução a \(W_q=2A_q-qI\) obtém-se
\[
\overline{W_q}=2\overline{A_q}-qI=-W_q.
\]
Logo \(W_q\) é anti-invariante. Como \(\mathcal{A}_q=\operatorname{span}_{\mathbb{Q}}\{I,W_q\}\), a involução induz a decomposição em eigenspaces
\[
\mathcal{A}_q=\mathcal{A}_q^+\oplus\mathcal{A}_q^-,
\]
onde
\[
\boxed{
\mathcal{A}_q^+=\mathbb{Q}I,
\qquad
\mathcal{A}_q^-=\mathbb{Q}W_q.
}
\]
A parte \(+\) é fixa e a parte \(-\) é anti-fixa. Esta é uma estrutura de paridade interna da álgebra, consequência direta das relações já estabelecidas; não constitui uma dualidade importada de estruturas externas.

A simetria de parâmetro \(q\mapsto-q\) é compatível com a involução: tem-se
\[
P_{-q}(t)=P_q(-t)
\]
e, no nível dos autovalores,
\[
\alpha_{-q}^+=\frac1{\alpha_q^+}.
\]

\section{Discriminante \(\Delta_q=q^2+4\)}

\begin{definicao}
O discriminante associado a \(q\) é
\[
\Delta_q:=q^2+4.
\]
\end{definicao}

Observa-se que \(\Delta_q\ge 4>0\) para todo \(q\in\mathbb{Q}\), logo
\(\Delta_q\neq 0\). Em particular \(W_q\) é sempre invertível e
\[
W_q^{-1}=\frac{1}{\Delta_q}W_q.
\]

\section{Fibras split e parametrização completa}

\begin{proposicao}
As seguintes condições são equivalentes:
\begin{enumerate}[label=(\roman*)]
\item \(\Delta_q\) é um quadrado em \(\mathbb{Q}\);
\item \(P_q\) se fatore em dois fatores lineares distintos sobre \(\mathbb{Q}\);
\item \(\mathcal{A}_q\cong\mathbb{Q}\times\mathbb{Q}\);
\item os autovalores de \(A_q\) pertencem a \(\mathbb{Q}\).
\end{enumerate}
Caso contrário, \(P_q\) é irredutível e \(\mathcal{A}_q\) é um corpo
quadrático (isomorfo a \(\mathbb{Q}(\sqrt{\Delta_q})\)).
\end{proposicao}

\begin{teorema}[Parametrização das fibras split]
As soluções racionais da equação \(r^2=q^2+4\) são dadas por
\[
q=\frac{4-s^2}{2s},\qquad
r=\frac{4+s^2}{2s},
\qquad
s\in\mathbb{Q}^\times.
\]
Em particular, o conjunto dos parâmetros \(q\) para os quais a fibra é split
é precisamente
\[
\Bigl\{\frac{4-s^2}{2s}:s\in\mathbb{Q}^\times\Bigr\}.
\]
\end{teorema}

\begin{proof}
A equação \(r^2-q^2=4\) escreve-se \((r-q)(r+q)=4\). Pondo \(s=r-q\) e
\(t=r+q\), tem-se \(st=4\). Resolvendo o sistema linear obtém-se a
parametrização acima. A verificação direta confirma que toda escolha de
\(s\in\mathbb{Q}^\times\) produz uma solução, e que toda solução racional
surge dessa forma.
\end{proof}

Assim a classificação das fibras é completa:
\[
\mathcal{A}_q\cong
\begin{cases}
\mathbb{Q}\times\mathbb{Q}
& \text{se }\Delta_q\in(\mathbb{Q}^\times)^2,\\[1mm]
\mathbb{Q}(\sqrt{\Delta_q})
& \text{se }\Delta_q\notin(\mathbb{Q}^\times)^2.
\end{cases}
\]

\section{Extensão mínima \(K_q\) e diagonalização}

\begin{definicao}
A extensão mínima associada a \(q\) é
\[
K_q:=\mathbb{Q}(\sqrt{\Delta_q}).
\]
\end{definicao}

\begin{proposicao}
\begin{itemize}
\item Se \(\Delta_q\) é quadrado em \(\mathbb{Q}\), então \(K_q=\mathbb{Q}\).
\item Caso contrário, \(K_q/\mathbb{Q}\) é uma extensão quadrática.
\end{itemize}
\end{proposicao}

Os autovalores de \(A_q\) são
\[
\lambda_\pm=\frac{q\pm\sqrt{\Delta_q}}{2}.
\]
Eles satisfazem
\[
\lambda_++\lambda_-=q,\qquad\lambda_+\lambda_-=-1.
\]
Sobre \(K_q\), \(A_q\) é sempre diagonalizável (os autovalores são distintos
porque \(\Delta_q\neq 0\)).

Uma matriz de mudança de base explícita é
\[
P=\begin{pmatrix}\lambda_+&\lambda_-\\1&1\end{pmatrix},
\]
que satisfaz
\[
P^{-1}A_qP=\operatorname{diag}(\lambda_+,\lambda_-)
\]
sempre que \(\det P=\sqrt{\Delta_q}\neq 0\). Quando a fibra é split tem-se
\(P\in M_2(\mathbb{Q})\).

\section{Similaridade dos \(A_q\)}

\begin{proposicao}
\[
A_q\sim A_r\quad\Longleftrightarrow\quad q=r.
\]
\end{proposicao}

\begin{proof}
O traço é invariante por similaridade e \(\operatorname{tr}(A_q)=q\).
A recíproca é trivial.
\end{proof}

A classe de similaridade de \(A_q\) contém outras matrizes além de \(A_q\)
(por exemplo \(T_q\)). Assim
\[
q\neq A_q\neq[A_q]_{\mathrm{similaridade}}.
\]
A parametrização por \(q\) é mais fina do que a classificação abstrata das
álgebras: diferentes valores de \(q\) podem produzir álgebras
abstratamente isomorfas (por exemplo quando \(\Delta_q\) e \(\Delta_r\)
estão na mesma classe módulo quadrados), enquanto \(A_q\) e \(A_r\) só
são semelhantes quando \(q=r\).

\section{Forma \(B_q\) e isotropia racional}

A matriz \(W_q\) é simétrica. Define-se a forma bilinear associada
\[
B_q(u,v)=u^\top W_q v
=
q\,u_1v_1+2u_1v_2+2u_2v_1-q\,u_2v_2,
\]
onde \(u=(u_1,u_2)\) e \(v=(v_1,v_2)\). A forma quadrática associada é
\[
Q_q(u):=B_q(u,u)=q\,u_1^2+4u_1u_2-q\,u_2^2.
\]
Tem-se
\[
\operatorname{tr}W_q=0,\qquad\det W_q=-\Delta_q.
\]
Como \(\Delta_q>0\), a forma é não-degenerada. Sobre \(\mathbb{R}\) toda
forma bilineares simétrica não-degenerada de dimensão \(2\) com
determinante negativo possui assinatura \((1,1)\). Logo todas as formas
\(B_q\) (e as formas quadráticas \(Q_q\)) são isométricas sobre \(\mathbb{R}\).

\subsection{Isotropia sobre \(\mathbb{Q}\)}

Um vetor \(u\neq0\) é isotrópico quando \(Q_q(u)=B_q(u,u)=0\). Para
\(q=0\) a forma \(Q_0(u)=4u_1u_2\) é claramente isotrópica. Para
\(q\neq0\), pondo \(t=u_1/u_2\) (com \(u_2\neq0\)), a equação torna-se
\(qt^2+4t-q=0\), cujo discriminante é \(4\Delta_q\). Logo existe solução
racional se e somente se \(\Delta_q\) é um quadrado em \(\mathbb{Q}\). Em
resumo:
\[
\boxed{
Q_q\text{ é isotrópica sobre }\mathbb{Q}
\quad\Longleftrightarrow\quad
\Delta_q\in(\mathbb{Q}^\times)^2.
}
\]
Assim o mesmo discriminante que classifica as fibras \(\mathcal{A}_q\)
controla também a isotropia racional da forma quadrática \(Q_q\).

\subsection{Regime split}

Quando \(\Delta_q=r^2\) com \(r\in\mathbb{Q}^\times\) e \(q\neq0\), as raízes
\[
t_\pm=\frac{-2\pm r}{q}
\]
fornecem vetores isotrópicos. A matriz
\[
P=\begin{pmatrix}t_+&t_-\\1&1\end{pmatrix}
\]
pertence a \(\mathrm{GL}_2(\mathbb{Q})\) (\(\det P=2r/q\)) e satisfaz
\[
P^\top W_q P
=
-\frac{2\Delta_q}{q}
\begin{pmatrix}0&1\\1&0\end{pmatrix}.
\]
O fator de escala \(\lambda=-2\Delta_q/q\neq0\) não altera o grupo
ortogonal: \(O(\lambda H;\mathbb{Q})=O(H;\mathbb{Q})\). Portanto
\[
\boxed{
O(B_q;\mathbb{Q})\cong O(1,1;\mathbb{Q})
}
\]
sempre que \(\Delta_q\) é um quadrado racional. (O caso \(q=0\) já havia
sido verificado por mudança de base independente.)

\subsection{Regime não-split}

Quando \(\Delta_q\notin(\mathbb{Q}^\times)^2\) a forma quadrática \(Q_q\) é
anisotrópica sobre \(\mathbb{Q}\) (apesar de indefinida sobre \(\mathbb{R}\)).
A descrição explícita do grupo \(O(B_q;\mathbb{Q})\) em termos de unidades
do corpo \(\mathbb{Q}(\sqrt{\Delta_q})\) não é desenvolvida aqui.

\subsection{Observação metodológica}

O discriminante \(\Delta_q\) aparece em três construções distintas
(\(\mathcal{A}_q\), \(B_q\), \(O(B_q;\mathbb{Q})\)). Isso não autoriza
identificá-las:
\[
\boxed{
\text{álgebra }\mathcal{A}_q
\;\neq\;
\text{forma }B_q
\;\neq\;
O(B_q;\mathbb{Q}).
}
\]
A forma \(B_q\) continua sendo uma construção adicional; não constitui,
por si só, uma geometria canônica associada a \(q\). Em particular
\[
\boxed{
\text{dependência algébrica em }q
\;\neq\;
\text{dependência geométrica real em }q.
}
\]

\section{Estruturas externas: ramos de Gentil e formas multiplicativas}

As operações por partes reportadas por Gentil constituem uma estrutura
\textbf{externa} em análise. Elas não fazem parte do núcleo algébrico da
família \(\mathcal{A}_q\) e não são identificadas com ela neste texto.
A postura é de confronto formal, não de assimilação.

A regra reportada é
\[
\mu(a,b;c,d)=(ac+\varepsilon(c)bd,\,ad+bc),
\]
onde
\[
\varepsilon(c)=\begin{cases}
-1 & \text{se }c\ge0,\\
+1 & \text{se }c<0.
\end{cases}
\]
Isso determina dois ramos constantes:
\begin{align*}
\mu_-(x,y)&=(ac-bd,\,ad+bc),\\
\mu_+(x,y)&=(ac+bd,\,ad+bc).
\end{align*}

Para \(x=(a,b)\) define-se inicialmente a norma euclidiana
\[
N(x)=a^2+b^2.
\]
No ramo negativo obtém-se a identidade clássica
\[
N(\mu_-(x,y))=N(x)N(y).
\]
Assim \(\mu_-\) coincide com a multiplicação complexa usual em coordenadas
reais e preserva a forma positiva \(a^2+b^2\).

No ramo positivo a situação é diferente:
\[
N(\mu_+(x,y))=(a^2+b^2)(c^2+d^2)+4abcd,
\]
logo a norma euclidiana \textbf{não} é multiplicativa em geral. Em
compensação, a forma indefinida
\[
Q_+(a,b)=a^2-b^2
\]
satisfaz
\[
Q_+(\mu_+(x,y))=Q_+(x)Q_+(y).
\]
Resumindo:
\[
\begin{array}{c|c|c}
\text{ramo} & \text{produto} & \text{forma multiplicativa}\\
\hline
- & (ac-bd,\,ad+bc) & a^2+b^2\\
+ & (ac+bd,\,ad+bc) & a^2-b^2
\end{array}
\]

Estas são duas multiplicatividades de formas quadráticas distintas,
correspondentes respectivamente às álgebras \(\mathbb{C}\) e
\(\mathbb{R}\times\mathbb{R}\) (ou, em linguagem de polinômios,
\(t^2+1\) e \(t^2-1\)).

A preservação de uma forma quadrática é uma propriedade algébrica
verificável da operação; não constitui, por si só, uma identificação
entre as estruturas.

\subsection{Classificação algébrica dos ramos fixos}

Os ramos \(\mu_+\) e \(\mu_-\) são tratados aqui como álgebras separadas
(produto constante, sem escolha de sinal dependente do estado). Os
resultados abaixo referem-se exclusivamente a esses ramos fixos.

\paragraph{Ramo \(\mu_+\): isomorfismo explícito.}
Com a adição componente a componente, \((\mathbb{Q}^2,\mu_+)\) é uma
\(\mathbb{Q}\)-álgebra comutativa de dimensão 2. Defina
\[
\Phi_+:\mathbb{Q}^2\to\mathcal{A}_0,
\qquad
\Phi_+(a,b)=aI+bA_0.
\]
Como \(A_0^2=I\) (pois \(P_0(t)=t^2-1\)) e
\[
\mu_+((a,b),(c,d))=(ac+bd,\,ad+bc),
\]
tem-se
\begin{align*}
\Phi_+(\mu_+(x,y))
&=\Phi_+(ac+bd,\,ad+bc)
=(ac+bd)I+(ad+bc)A_0
\\
&=(aI+bA_0)(cI+dA_0)
=\Phi_+(x)\Phi_+(y).
\end{align*}
O mapa é linear e leva a base canônica de \(\mathbb{Q}^2\) na base
\(\{I,A_0\}\) de \(\mathcal{A}_0\), logo é bijetivo. Portanto
\[
\boxed{(\mathbb{Q}^2,\mu_+)\cong\mathcal{A}_0}.
\]
(DEMONSTRADO)

\paragraph{Ramo \(\mu_-\): obstrução para toda a família.}
A álgebra \((\mathbb{Q}^2,\mu_-)\) é isomorfa a \(\mathbb{Q}[i]\) com
\(i^2=-1\) (a multiplicação complexa usual). O polinômio mínimo de \(i\)
é \(t^2+1\), de discriminante \(-4<0\).

Por outro lado, para todo \(q\in\mathbb{Q}\),
\[
\mathcal{A}_q\cong\mathbb{Q}[t]/(t^2-qt-1)
\]
tem discriminante \(\Delta_q=q^2+4>0\). Duas \(\mathbb{Q}\)-álgebras
quadráticas isomorfas possuem discriminantes que diferem por um
quadrado em \(\mathbb{Q}^\times\). Como um discriminante positivo nunca
é um quadrado racional vezes um discriminante negativo, nenhuma fibra
\(\mathcal{A}_q\) pode ser isomorfa a \(\mathbb{Q}[i]\). Logo
\[
\boxed{(\mathbb{Q}^2,\mu_-)\not\cong\mathcal{A}_q\quad\forall q\in\mathbb{Q}}.
\]
(DEMONSTRADO)

\paragraph{Alcance da conclusão.}
Os resultados acima valem para os ramos \emph{fixos} \(\mu_+\) e
\(\mu_-\), tratados como álgebras separadas. Eles \textbf{não}
demonstram que a operação original de Gentil (caso esta seja uma
operação global definida por escolha de sinal dependente do estado)
seja isomorfa a \(\mathcal{A}_0\), nem estabelecem qualquer relação com
estruturas tridimensionais. Em particular, não se afirma dualidade,
identificação global ou “caso limite” entre Gentil e a família
\(\{\mathcal{A}_q\}\).

Uma comparação adicional legítima (ainda em nível de observação) é a
seguinte: a forma \(B_q\) associada a \(W_q\) tem sempre assinatura real
\((1,1)\), enquanto \(Q_+\) é a forma indefinida padrão de dimensão 2.
No caso particular \(q=0\) tem-se
\[
W_0=\begin{pmatrix}0&2\\2&0\end{pmatrix},
\]
logo
\[
B_0\begin{pmatrix}u\\v\end{pmatrix}=4uv.
\]
A identidade
\[
4uv=(u+v)^2-(u-v)^2
\]
mostra que, sob a mudança de variáveis
\[
\begin{pmatrix}a\\b\end{pmatrix}
=
\begin{pmatrix}1&1\\1&-1\end{pmatrix}
\begin{pmatrix}u\\v\end{pmatrix},
\]
obtém-se
\[
B_0\begin{pmatrix}u\\v\end{pmatrix}=a^2-b^2.
\]
A matriz da mudança tem determinante \(-2\neq0\), portanto é invertível
sobre \(\mathbb{Q}\). A equivalência racional é assim explícita e
verificável. Essa equivalência pontual não se estende automaticamente
a uma identificação das estruturas.

\section{Família global e distinção entre parâmetro e tipo de álgebra}

Os objetos matemáticos distintos são:
\begin{itemize}
\item o espaço de parâmetros \(\mathbb{Q}\) (ou \(\mathcal{P}_{\mathbb{Q}}\));
\item a família de polinômios \(\{P_q:q\in\mathbb{Q}\}\);
\item a família de matrizes \(\{A_q:q\in\mathbb{Q}\}\);
\item a família de álgebras \(\{\mathcal{A}_q:q\in\mathbb{Q}\}\).
\end{itemize}

A união disjunta
\[
\bigsqcup_{q\in\mathbb{Q}}\mathcal{A}_q
\]
é bem-definida como conjunto. Cada fibra carrega sua própria estrutura de
\(\mathbb{Q}\)-álgebra de dimensão \(2\). Não existe, no estado atual,
estrutura de álgebra (nem de espaço vetorial) no total que misture
diferentes parâmetros.

O objeto mais fiel aos dados é a
\[
\boxed{
\text{família de álgebras quadráticas parametrizada por }\mathbb{Q}.
}
\]
Não se introduz estrutura de moduli, fibrado ou variedade sem justificativa
adicional.

\section{Relação com a implementação Tiffany}

A aritmética racional necessária para parametrizar \(q\) possui realização
computacional no Tiffany (tipo \(\mathtt{Qz}\), normalização por sinal do
denominador e redução por máximo divisor comum). A identidade
\[
A_q^2=qA_q+I
\]
e a relação \(W_q^2=\Delta_q I\) foram verificadas experimentalmente para
inteiros e simuladas para racionais.

A composição computacional direta \(q\mapsto A_q\) com \(q\in\mathbb{Q}\)
ainda requer uma camada matricial racional explícita ou uma generalização
do contrato matricial existente (atualmente centrado em parâmetros inteiros).

A assinatura prima
\[
\Sigma_{\mathrm{rad}}(x)=(\nu_p(x))_{p\in\mathbb{P}}
\]
não possui implementação computacional correspondente a uma aplicação
\[
\Sigma_{\mathrm{rad}}:\mathbb{Q}^*\to\mathbb{Z}^{(\mathbb{P})}
\]
no estado atual. Consequentemente a cadeia
\[
x\longmapsto\Sigma_{\mathrm{rad}}(x)\longmapsto q\longmapsto P_q
\]
não está realizada. Qualquer conexão futura deverá ser introduzida e
demonstrada separadamente.

O mecanismo JEV e a família \(\mathcal{P}_{\mathbb{Q}}\) permanecem
construções paralelas; não há aplicação demonstrada \(G\mapsto q\).

\section{Classificação epistemológica}

\begin{center}
\begin{tabular}{@{}p{0.55\textwidth}p{0.35\textwidth}@{}}
\toprule
Afirmação & Estatuto \\
\midrule
\(\mathcal{A}_q\cong\mathbb{Q}[t]/(P_q)\) & DEMONSTRADO \\
\(\mathcal{A}_q\cong\mathbb{Q}[w]/(w^2-\Delta_q)\) & DEMONSTRADO \\
Classificação split / corpo segundo \(\Delta_q\) & DEMONSTRADO \\
Parametrização racional completa das fibras split & DEMONSTRADO \\
\(A_q\sim A_r\Leftrightarrow q=r\) & DEMONSTRADO \\
Extensão mínima \(K_q=\mathbb{Q}(\sqrt{\Delta_q})\) & DEMONSTRADO \\
Diagonalização explícita sobre \(K_q\) & DEMONSTRADO \\
Involução \(\overline{A_q}=qI-A_q\), \(\overline{\overline{A_q}}=A_q\) & DEMONSTRADO \\
\(A_q\overline{A_q}=-I\) & DEMONSTRADO \\
Paridade \(\mathcal{A}_q=\mathbb{Q}I\oplus\mathbb{Q}W_q\) & DEMONSTRADO \\
\(\det W_q=-\Delta_q\), assinatura real \((1,1)\) & DEMONSTRADO \\
Isometria real de todas as \(B_q\) & DEMONSTRADO \\
\(Q_q\) isotrópica \(\Leftrightarrow\Delta_q\in(\mathbb{Q}^\times)^2\) & DEMONSTRADO \\
Matriz \(P\) e fator \(\lambda=-2\Delta_q/q\) no caso split & DEMONSTRADO \\
\(O(B_q;\mathbb{Q})\cong O(1,1;\mathbb{Q})\) quando \(\Delta_q\) é quadrado & DEMONSTRADO \\
Descrição de \(O(B_q;\mathbb{Q})\) no regime não-split & NÃO ESTABELECIDO \\
\(N(\mu_-)=N(x)N(y)\) (forma \(a^2+b^2\)) & DEMONSTRADO \\
\(Q_+(\mu_+)=Q_+(x)Q_+(y)\) (forma \(a^2-b^2\)) & DEMONSTRADO \\
\((\mathbb{Q}^2,\mu_+)\cong\mathcal{A}_0\) & DEMONSTRADO \\
\((\mathbb{Q}^2,\mu_-)\not\cong\mathcal{A}_q\) para todo \(q\in\mathbb{Q}\) & DEMONSTRADO \\
Identificação da operação global de Gentil com \(\mathcal{A}_q\) & NÃO ESTABELECIDO \\
Classificação completa de isometria racional de \(B_q\) & NÃO ESTABELECIDO \\
Estrutura geométrica canônica por fibra & NÃO ESTABELECIDO \\
Operações naturais entre parâmetros distintos & NÃO ESTABELECIDO \\
Cadeia \(\Sigma_{\mathrm{rad}}\to q\) & NÃO ESTABELECIDO \\
\bottomrule
\end{tabular}
\end{center}

\section{Conclusão}

A construção sustentada estabelece a cadeia matemática
\[
\boxed{
q\in\mathbb{Q}
\longmapsto
P_q(\lambda)=\lambda^2-q\lambda-1
\longmapsto
A_q
\longmapsto
\mathcal{A}_q
\cong
\mathbb{Q}[t]/(P_q)
\cong
\mathbb{Q}[w]/(w^2-\Delta_q)
}
\]
com \(\Delta_q=q^2+4\), e a dicotomia
\[
\Delta_q\in(\mathbb{Q}^\times)^2
\quad\Longleftrightarrow\quad
\mathcal{A}_q\cong\mathbb{Q}\times\mathbb{Q},
\]
\[
\Delta_q\notin(\mathbb{Q}^\times)^2
\quad\Longleftrightarrow\quad
\mathcal{A}_q\cong\mathbb{Q}(\sqrt{\Delta_q}).
\]

A parametrização por \(q\) é mais fina que a classificação abstrata das
álgebras. O mesmo discriminante controla independentemente a natureza
das fibras e a isotropia das formas quadráticas:
\[
\Delta_q\in(\mathbb{Q}^\times)^2
\quad\Longleftrightarrow\quad
\mathcal{A}_q\cong\mathbb{Q}\times\mathbb{Q}
\quad\text{e}\quad
Q_q\text{ isotrópica},
\]
com a consequência
\[
O(B_q;\mathbb{Q})\cong O(1,1;\mathbb{Q})
\]
no regime split. No regime não-split a forma quadrática \(Q_q\) é
anisotrópica sobre \(\mathbb{Q}\) e a descrição do grupo ortogonal não é
desenvolvida. A coincidência do invariante não identifica os objetos:
\[
\mathcal{A}_q\;\neq\;B_q\;\neq\;O(B_q;\mathbb{Q}).
\]

A álgebra \(\mathcal{A}_q\) possui uma involução interna
\[
\overline{A_q}=qI-A_q
\]
que induz a paridade
\[
\mathcal{A}_q=\mathbb{Q}I\oplus\mathbb{Q}W_q.
\]
Esta é estrutura interna, consequência direta das relações já estabelecidas.

Os ramos fixos de Gentil foram classificados algebricamente:
\[
(\mathbb{Q}^2,\mu_+)\cong\mathcal{A}_0,
\qquad
(\mathbb{Q}^2,\mu_-)\not\cong\mathcal{A}_q\quad\forall q\in\mathbb{Q}.
\]
Estes resultados referem-se exclusivamente aos ramos constantes;
não se estendem automaticamente à operação global (com escolha de
sinal dependente do estado) nem a estruturas tridimensionais.

O objeto matemático central é a família de álgebras quadráticas
parametrizada por \(\mathbb{Q}\). Nenhuma geometria canônica, nenhuma
operação entre parâmetros distintos e nenhuma estrutura global mais forte
foram introduzidas.

A ordem metodológica preservada é
\[
\boxed{
\text{realização}
\rightarrow
\text{representação}
\rightarrow
\text{estrutura algébrica}
\rightarrow
\text{simetria interna (involução/paridade)}
\rightarrow
\text{classificação aritmética das fibras}
\rightarrow
\text{estruturas externas em confronto}
\rightarrow
\text{possível construção geométrica}.
}
\]

\begin{thebibliography}{9}

\bibitem{geometria-observador-m}
Auditoria do Tiffany --- geometria observacional e realização aritmética.
\texttt{geometria\_observador\_md}.
Seções relevantes: estruturas reais, cadeia de realização,
extensão \(\mathbb{Z}\!\to\!\mathbb{Q}\), construção \(P_q\),
matriz \(T_q\), estruturas geométricas, assinatura prima, JEV,
classificação final e auditoria cirúrgica \(\mathbb{Q}\!\to\!P_q\).

\end{thebibliography}

\end{document}
